% Advances in Water Resources submission.
%
% elsarticle option notes:
%   preprint    single column, single spaced  <- used here
%   review      single column, double spaced, line numbers (add for review copy)
%   3p,twocolumn  the printed double-column layout -- NOT wanted here
%   authoryear  REQUIRED for (Author year) citations. elsarticle defaults to
%               numbered mode, so without it \citep renders as [5] no matter which
%               .bst is loaded, and elsarticle-harv silently produces a numbered
%               reference list.
%
% EMS does not require a specific reference format at submission; any consistent
% style is accepted and the journal style is applied at proof stage. The journal's
% published in-text form is author-date, so elsarticle-harv is used.
\documentclass[preprint,authoryear,11pt]{elsarticle}

\usepackage[T1]{fontenc}
% elsarticle's preprint option sets a 360 pt (5.0 in) measure, which leaves 1.75 in
% of margin each side on letter paper. The journal does not specify a text width --
% Elsevier applies its own layout at proof stage and the submitted PDF is a reading
% copy -- so the measure is ours to choose. Widened to 6.5 in. Keep the margin at or
% above 1 in: \widefig draws at 190 mm (7.48 in), and the overhang it relies on
% disappears once the text block reaches that width.
\usepackage[margin=1in]{geometry}
\usepackage{amsmath}
\usepackage[strings]{underscore}
\usepackage{booktabs}
\usepackage{array}
\usepackage{siunitx}
% Quantities are reported in SI throughout, with no exceptions. The sewer model is
% metric (gp_sewer.inp declares FLOW_UNITS CMS), so its lengths are given as it
% records them; MODFLOW is in feet internally, which the exchange formulations
% convert rather than report. The sewer-infiltration acceptance band is published
% in gallons per day per inch of diameter per mile and is converted to liters per
% day per millimeter of diameter per kilometer, a factor of 0.09260; the customary
% form is named once, where the band is introduced, so that a reader who knows the
% standard can find it. Units of that shape are spelled out in prose, so no unit
% needs declaring here.
% siunitx does not group digits by default, so a four-digit \SI would print "9550"
% in the same sentence as a hand-typed "6,491". Group at four digits so the two
% forms agree; the missing comma is invisible in the source.
%
% group-digits=integer is not optional here. Grouping applies to BOTH sides of the
% decimal marker by default, so the four digits of \SI{0.0655}{\per\day} were being
% grouped as well and the leakance typeset as "0.065,5" -- which reads as two
% numbers rather than one. Restricting it to the integer part leaves the thousands
% separator the comment above wants and lets a decimal run as long as it needs to.
\sisetup{group-separator={,}, group-minimum-digits=4, group-digits=integer}
\usepackage{graphicx}
% elsarticle 3.5 has no native ORCID support; orcidlink renders the iD as a
% hyperlinked icon after the author name.
\usepackage{orcidlink}
% elsarticle loads natbib itself -- do not load it again or the options clash.
\setcitestyle{aysep={}}          % (Author year), no comma
\usepackage{xcolor}
\usepackage{hyperref}
% Long repository URLs in the software statement cannot break in a 5 in measure and
% overflow the margin; xurl allows a break at any character.
\usepackage{xurl}
\hypersetup{colorlinks=true, linkcolor={red!50!black},
  citecolor={blue!50!black}, urlcolor={blue!80!black}}
\newcolumntype{L}[1]{>{\raggedright\arraybackslash}p{#1}}

% Full-width figures are drawn at 190 mm, which is the journal's full-page artwork
% width and the width the figure scripts author to. The preprint text block is only
% 360 pt (5.0 in), so \includegraphics[width=\textwidth] would scale them to 67% and
% render 6.5 pt annotations at 4.3 pt -- artwork "where text is disproportionately
% small", which the guide lists among the things not to submit. Prose keeps the
% narrower measure, which is easier to read, and the figures overhang it
% symmetrically; on letter paper that still leaves 0.51 in of margin each side.
% The box is declared zero width so the overhang raises no overfull warning.
\newlength{\figfull}
\setlength{\figfull}{190mm}
\newcommand{\widefig}[1]{\centerline{\makebox[0pt][c]{%
  \includegraphics[width=\figfull]{#1}}}}

% Single-column artwork is drawn at 90 mm, the journal's one-column width, and is
% placed 1:1 for the reason \widefig is: the figure scripts author to a stated width
% and scaling one defeats the font sizes chosen for it. A figure that carries a
% single statement is drawn at this width rather than stretched to fill the measure.
% It is centered in the text block, and needs no overhang.
\newlength{\figcol}
\setlength{\figcol}{90mm}
\newcommand{\onecolfig}[1]{\centerline{\includegraphics[width=\figcol]{#1}}}

\journal{Advances in Water Resources}

\begin{document}

\begin{frontmatter}

% Title note: "library interfaces" rather than "the Basic Model Interface".
% MODFLOW 6 and D-Flow FM implement BMI; SWMM does not, and is driven through the
% SWMM toolkit API, which supplies the same capabilities under different names.
% Naming BMI in the title would misstate one of the three. BMI is carried in the
% abstract and keywords instead, where it can be qualified.
%
% The title states the finding rather than the method, which is what the AdWR scope
% statement asks for: a paper reading as an application of existing knowledge is
% explicitly out of scope there. It claims nothing beyond what the conclusions
% already say in those words -- 12-hour exchange carries more head error than
% 24-hour, because both alias the M2 constituent onto the same 14.8-day beat -- and
% that is a property of aliasing rather than of this site. The method title it
% replaces was written for Environmental Modelling & Software:
%   Dynamic coupling of coastal hydrodynamic, groundwater flow, and sewer network
%   models using library interfaces
\title{A shorter coupling interval is not a safer one: tidal aliasing in coupled
coastal hydrodynamic, groundwater, and sewer simulations}

% Name spellings and affiliations below are taken from each author's ORCID record
% (pub.orcid.org/v3.0/<id>/person and /employments), not from correspondence, so
% they should be confirmed by each author before submission. Office assignments
% for the New York Water Science Center authors follow their ORCID employment
% city: Coram for Bayraktar, Troy for Jahn.
% Affiliation labels follow FIRST APPEARANCE in the author list, and the
% \affiliation declarations are given in that same order. elsarticle letters them in
% declaration order, so any other arrangement prints the superscripts out of
% sequence: with Prigiobbe added at inst5 in third position they read a, b, e, c, d,
% b. Inserting an author means renumbering from that author onward.
% The guide requires a corresponding author to be marked. The address below is a
% PLACEHOLDER and must be replaced before submission -- it is the only detail here
% not taken from an ORCID record.
% The first author carries two affiliations because the work spans an employer
% change: it began at the USGS Integrated Modeling and Prediction Division
% (ORCID employment 2021 to 2025-03, and the commits under jdhughes@usgs.gov run
% 2024-11 to 2025-02) and continued at INTERA. Affiliation records where work was
% performed, so dropping either one would misattribute part of it.
\author[inst1,inst2]{Joseph D. Hughes\corref{cor1}\,\orcidlink{0000-0003-1311-2354}}
\cortext[cor1]{Corresponding author. \textit{Email address:} jdhughes@intera.com}
\author[inst3]{Liv Herdman\,\orcidlink{0000-0002-5444-6441}}
\author[inst4]{Valentina Prigiobbe\,\orcidlink{0000-0001-6603-4295}}
\author[inst5]{Martijn J. Russcher\,\orcidlink{0000-0001-8799-6514}}
\author[inst6]{Banu Bayraktar\,\orcidlink{0000-0003-3612-6767}}
\author[inst3]{Kalle L. Jahn\,\orcidlink{0000-0002-4976-0137}}

% The Integrated Modeling and Prediction Division duty station and the INTERA
% office are the same Chicago address, so inst1 and inst2 repeat it. That is not a
% copy-and-paste error: the first author worked from this address at both
% employers, which is also why the ORCID employment record gives Chicago for the
% division. Division, city and state are from that record.
\affiliation[inst1]{organization={U.S. Geological Survey, Integrated Modeling and Prediction Division}, addressline={927 W Belle Plaine Ave}, city={Chicago}, state={IL}, postcode={60613}, country={USA}}

\affiliation[inst2]{organization={INTERA Incorporated}, addressline={927 W Belle Plaine Ave}, city={Chicago}, state={IL}, postcode={60613}, country={USA}}

\affiliation[inst3]{organization={U.S. Geological Survey, New York Water Science Center}, addressline={425 Jordan Road}, city={Troy}, state={NY}, postcode={12180}, country={USA}}

\affiliation[inst4]{organization={Department of Geosciences, University of Padua}, addressline={Via Gradenigo 6}, city={Padua}, postcode={35131}, country={Italy}}

\affiliation[inst5]{organization={Deltares}, addressline={Boussinesqweg 1}, city={Delft}, postcode={2629 HV}, country={The Netherlands}}

\affiliation[inst6]{organization={U.S. Geological Survey, New York Water Science Center}, addressline={2045 Route 112, Building 4}, city={Coram}, state={NY}, postcode={11727}, country={USA}}

% SUBMISSION REQUIREMENTS -- Advances in Water Resources (ISSN 0309-1708).
% Read off the guide for authors, 2026-08-20, on retargeting from Environmental
% Modelling & Software. Where the two journals differ the EMS figure has been
% replaced rather than adapted, and where they agree that is stated, so the next
% person does not have to re-check.
%
% SCOPE -- the paragraph to read before any of the formatting. AdWR asks for
% "fundamental scientific advances in the understanding of water resources
% systems" and rules out "case studies that do not attempt to reach broader
% conclusions" and "applications of existing knowledge that do not advance
% fundamental understanding of hydrological processes". Greenport is therefore
% presented as the problem that exercises all three couplings at once and never as
% a site characterization, and the claim the paper makes is the general one: how
% the coupling interval, and the representation of the tidal boundary over it, set
% the error of the coupled solution. Guard that framing in revision.
%
% IN THIS FILE
%   Abstract     250 words maximum, up from the 150 EMS allowed, which is why it
%                is no longer written against a word budget. Must stand alone,
%                avoid references, and define any uncommon abbreviation.
%   Keywords     1 to 7; seven here. Avoid keywords joined by "and" or "of", and
%                abbreviations unless firmly established.
%   Sections     NUMBERED -- which OVERRIDES the unnumbered house style, for this
%                submission only -- subsections as 1.1 and then 1.1.1, cross-referenced by
%                number rather than as "the text". The abstract is excluded from
%                the numbering; the back matter is unnumbered.
%   Equations    Numbered consecutively in order of first reference. Within an
%                appendix, Eq. (A.1), Table A.1, Fig. A.1.
%   Generative AI declaration: its own section before the reference list, which is
%                where it sits. The guide's section title is used verbatim.
%   Corresponding author must be marked. See \cortext below.
%
% SEPARATE FILES AT SUBMISSION -- none of these are part of this .tex
%   Highlights         REQUIRED. 3 to 5 bullets, 85 characters maximum each,
%                      including spaces, as a separate editable file with
%                      "highlights" in the name. Identical to the EMS requirement,
%                      so the existing file carries over unchanged.
%   Graphical abstract ENCOURAGED, not required. This is the one relaxation over
%                      EMS, and the specification is the same one:
%                      531 x 1328 px (h x w) or proportionally more, readable at
%                      5 x 13 cm, as TIFF, EPS, PDF or MS Office. So
%                      scripts/make_graphical_abstract.py needs no change, and the
%                      graphical abstract can now be dropped rather than fixed if
%                      it ever becomes a nuisance.
%   Figures            Separate files named Figure_1, Figure_2, ... Vector art as
%                      EPS or PDF with fonts embedded, which is what the scripts in
%                      scripts/ already produce.
%   Declarations       Competing interests, through Elsevier's declarations tool,
%                      as a .doc/.docx upload.
%
% STILL OUTSTANDING -- see the back matter for the sections themselves
%   CRediT, the funding statement, and the data deposit need author input before
%   submission. Funding is a required disclosure and the manuscript carried no such
%   section under EMS; there is now a visible placeholder in the back matter.
%   Research data: AdWR is Option C, exactly as EMS was, so the deposit must be
%   made, cited, and linked in the article, or a statement given explaining why it
%   cannot be shared. A data statement is required at submission either way.
%
% BACK MATTER, in order: data and code availability, CRediT, the generative-AI
% declaration, competing interests, funding, disclaimer, references. The guide's one
% ordering rule concerns acknowledgements, which this manuscript does not carry.
%
%   https://www.sciencedirect.com/journal/advances-in-water-resources/publish/guide-for-authors
\begin{abstract}
A coupled simulation exchanges state at an interval that is usually chosen from run time rather than from the processes being resolved. Three unmodified simulators are loaded as libraries into one process and exchange state in memory, which leaves that interval free to vary; returning sewer discharge to the hydrodynamic model requires a transfer within a model time step that the Basic Model Interface does not define. It is exercised where a tidally modulated water table infiltrates a combined sewer, across three grids and intervals from 15 minutes to one day. The tidal boundary must represent the interval, not the instant: sampling it costs an order of magnitude in head at daily exchange and delivers a quarter more water to the aquifer. The error is not monotonic, being larger at 12 hours than at 24, following the period onto which the dominant tidal constituent aliases, so a shorter interval cannot be assumed to be a safer one. Below the Nyquist limit the two representations differ by amounts too small to matter, and which of them is the more accurate depends on the instant the simulation begins. Net infiltration to the sewer answers to the hydrodynamic grid instead, varying more than twofold across the three and crossing the leakage a new sewer is permitted. Exchanging often is inexpensive, coupling 96 times more frequently costing under 20 percent more run time on the finest grid, so the interval is better chosen from the tidal constituents present than from the run-time budget.
\end{abstract}

\begin{keyword}
coupling interval \sep
tidal aliasing \sep
groundwater--surface water interaction \sep
compound flooding \sep
sewer infiltration \sep
coastal aquifer \sep
model coupling \sep
Basic Model Interface
\end{keyword}

\end{frontmatter}

\section{Introduction}

Processes that cross the boundary between the coastal ocean, the subsurface, and engineered infrastructure are rarely represented by a single code. Each domain has a mature simulator of its own, developed independently, with its own governing equations, spatial discretization, model time step, and means of external control. Representing the coupled process therefore means making those codes exchange state while they run.

Coupling of this kind has generally been achieved in one of two ways. The first adds the second process to an existing code, which yields a tightly integrated solution but confines it to the domains the host code implements; the coupled surface water routing added to MODFLOW is an example \citep{hughes2014swr}. The second couples separate codes externally, and has been applied to urban drainage and groundwater through SWMM and MODFLOW \citep{zhang2020swmmmodflow, li2022urbansaturated, tajari2026compound}, to large-scale hydrology and groundwater \citep{guillaumot2022cwatm}, and to the coastal zone by coupling a hydrodynamic model to a groundwater model \citep{jin2025gwocean}. In that last the exchange interval is the 10-minute output cadence of the hydrodynamic model, against a 10-second hydrodynamic time step, chosen so that the two models stay synchronized; its effect on the coupled solution is not reported.

Standardized control interfaces make the second route practical without the cost that has historically accompanied it. Simulators that provide a library interface, such as the Basic Model Interface \citep{hutton2020bmi} or an application programming interface built for the same purpose \citep{hughes2022modflowapi}, can be loaded into a single process and advanced under external control, with state passed as arrays in memory. The coupling interval is then set by the numerics rather than by any intermediate file format, and a self-consistent coupled solution is obtained from one simulation.

A simulator need not implement the Basic Model Interface to be coupled. The only requirement is that the code be callable as a shared library and provide four capabilities. It must initialize a simulation from its own input or through memory, advance it by a time step, finalize it, and read and write named state. The Basic Model Interface codifies those four capabilities. Where a simulator implements it, the driver calls them by their standard names; where a simulator does not, an interface offering the same four under different names serves equally well. Of the three codes coupled here, two implement the Basic Model Interface and the third does not, and the distinction turns out to matter only to the names of the calls.

The decisive property is that none of the coupled codes is modified. Each is used as distributed: the coupling lives entirely in the controlling program, which advances each simulator through its documented interface and moves state between them. No source is forked, no process is reimplemented inside a host code, and each simulator continues to be developed, tested, and released by the group that maintains it. A coupling built this way survives a new release of any component, and the same controlling program can be applied to a different combination of simulators. \cite{jin2025gwocean} exploit the same property in coupling a hydrodynamic model to MODFLOW~6, adjusting the general-head boundary through the interface rather than by altering the code.

What such an interface provides, however, is not always sufficient on its own. A quantity may be readable but not writable in a way the solution uses, and where within the model time step a value is written can determine whether it survives to be used at all. The Basic Model Interface defines initialization, advancement in time, and finalization; it does not define access to model state part way through a time step. The MODFLOW application programming interface extends the Basic Model Interface precisely because the ability to modify variables within a time step falls outside it \citep{hughes2022modflowapi}. Any coupling that must inject state between two operations of a single time step therefore depends on something beyond the standard, and a simulator offering only the standard interface cannot support that class of coupling at all. Both simulators driven here provide such an extension, by different routes, and the coupling described below depends on it.

This study has two objectives. The first is to establish how the coupled solution depends on the interval at which state is exchanged and on how the tidal boundary is represented over that interval, both of which a coupling in memory leaves free and neither of which is ordinarily chosen from the tide. The second is to describe the approach that makes them measurable: coupling a coastal hydrodynamic model, a groundwater flow model, and a sewer network model in memory through their library interfaces, and setting out the exchange formulations between each pair. Returning the sewer model outfall discharge to the hydrodynamic model as a sewer discharge tracer source places the strongest requirement on the interface: the transfer must occur between the external forcing update and the flow computation, a window the combined time-stepping call such interfaces usually present does not provide, and the calls that requirement implies are set out in \ref{app:interface}.

The approach is demonstrated on a problem that exercises all three couplings at once. This work is a component of a regional assessment of compound flood hazard from the combined effects of coastal, stormwater, and groundwater emergence flooding in New York City, Long Island, and the Long Island Sound watersheds \citep{usgs2021compound}. That assessment ranks the hazard from rainfall, from coastal storm surge, and from shallow or emerging groundwater on a 900-meter grid \citep{usgs2026lismap}; resolving how those drivers interact at the scale a sewershed drains asks instead for a coupled model of the coastal, groundwater, and infrastructure systems together, which is what is described here. Of the three drivers it is groundwater emergence, a shallow water table intersecting the land surface or critical infrastructure, that bears most directly on the sewer.

Compound flooding is more often modeled by combining the drivers that arrive at the surface. A reduced-physics solver carries fluvial, pluvial, tidal, wind- and wave-driven processes in a single code at a cost low enough for large domains \citep{leijnse2021compound}, and a tightly coupled one- and two-dimensional model resolves the drainage network and the flood defenses of an urban site \citep{tang2025compound}. Surface and subsurface flows, where they are coupled at all in that literature, are commonly coupled one way, and infiltration into the pipe is set aside: the second of those studies neglects it in its one-dimensional domain and names groundwater among the processes it leaves for later \citep{tang2025compound}. It is that infiltration, arriving from below, which is represented here. The contribution of groundwater to coastal flooding is routinely overlooked \citep{seidenfaden2025compound}, and where a compound flood framework has coupled it two ways to a tidal boundary the exchange was set at one hour to match the groundwater model's stress periods, with no test of that choice reported \citep{pena2022compound}. Tidal fluctuation drives a two-way exchange of water between the land, groundwater, and surface waters, a connection Earth system models still reduce to a one-way freshwater discharge \citep{ward2020coastal}.

Sewers in low-lying coastal communities are subject to groundwater infiltration that varies with the tide: where the water table stands above the pipe invert, groundwater enters the sewer through defects in the pipe wall and at joints, adding to the volume the network conveys. The water table is itself tidally forced, so the infiltration rate is modulated at tidal frequencies by a process external to the sewer network. That interaction, and its consequences for flooding and for water quality, has been studied in coastal urban settings through field measurement, statistical and machine-learning models, and scenarios of sea level rise \citep{liu2018gwsewer, su2019infiltration, liu2021predicting, su2022slr}.

The demonstration uses the sewer network at Greenport, New York, on the north fork of eastern Long Island, which discharges to Greenport Harbor and the surrounding waters of Long Island Sound and Peconic Bay. The network is a combined system carrying both stormwater runoff and dry weather flow. It is used here to evaluate the coupling rather than to characterize the site, and in particular to establish how frequently the exchange must occur for the coupled solution to be insensitive to the coupling interval.

\section{Model components}

% TODO: the D-Flow FM Long Island Sound grid has no citable reference. Searches of
% Herdman's published works, USGS publications and Crossref turned up nothing that
% documents it, and the nearest Long Island compound-flooding paper
% (Glas et al., NHESS 2026) is statistical rather than hydrodynamic. If the grid is
% documented in a data release or report, cite it here; otherwise it needs a
% sentence describing its provenance.
Three simulators are coupled. Coastal hydrodynamics and the transport of a conservative tracer of sewer discharge are simulated with D-Flow Flexible Mesh (D-Flow FM), which solves the depth-averaged shallow water equations on an unstructured grid using the scheme of \cite{kernkamp2011unstructured} and is documented in \cite{deltares2026dflowfm}. The advection terms of that scheme are integrated explicitly \citep{kernkamp2011unstructured}, so its computational time step is limited by the Courant condition and varies through a simulation. What is set instead is a longer user time step, at which external forcings are updated and at which a controlling program can reach the model; the coupling is written against that step and never against the computational one. Groundwater flow is simulated with MODFLOW~6 \citep{langevin2017modflow6}, using the Groundwater Flow (GWF) and Groundwater Transport (GWT) models with the Buoyancy (BUY) package so that density effects associated with the saltwater interface are represented. Flow in the combined sewer network is simulated with the Storm Water Management Model (SWMM) \citep{rossman2022swmm}.

% Scope statement. The preceding paragraph gives D-Flow FM a conservative tracer and
% MODFLOW 6 a saltwater interface in consecutive sentences, which invites the reader
% to assume salinity crosses the coupling. It does not, and saying so here is
% cheaper than leaving a reviewer to find it. Verified rather than asserted: all 68
% FlowFM.mdu files in dflow-fm/ set Salinity = 0 and Temperature = 0. The
% source-sink field claim is documented at common/dflow_bmi_structures.py:88-93,
% where change_in_salinity and change_in_temperature return NULL precisely because
% isalt and itemp are zero.
Salinity is not simulated by the hydrodynamic model. D-Flow FM is run with salinity and temperature transport disabled, so its solution is barotropic and the conservative sewer discharge tracer is the only constituent it transports. The saltwater interface represented by the MODFLOW~6 Buoyancy package is accordingly driven by concentrations set within the groundwater model rather than by any exchanged with the coastal model. Nothing in the approach precludes coupling the two transport solutions: a concentration would be carried by the same mapping matrices that carry water level, and the D-Flow FM source-sink structure already presents a change in salinity alongside the discharge, a field that resolves once salinity is active. That exchange was not required for the demonstration reported here, which concerns the aquifer head, the aquifer--sewer seepage, and the sewer discharge tracer, and it is left as an extension.

The three simulators are advanced within a single process and exchange state through their respective library interfaces. MODFLOW~6 is controlled through its application programming interface, which extends the Basic Model Interface \citep{hughes2022modflowapi}, and D-Flow FM through its Basic Model Interface implementation \citep{hutton2020bmi}. SWMM implements neither, and is controlled through the SWMM toolkit by way of PySWMM \citep{mcdonnell2020pyswmm}; that interface opens a simulation from a SWMM input file, advances it by a routing step, closes it, and reads and writes node and link state, which are the same capabilities the other two supply under Basic Model Interface names. The MODFLOW~6 model is constructed with FloPy \citep{hughes2023flopy}. All three simulators are used as distributed, without modification.

% VERSION STATEMENT -- must match whichever runs produce the reported results.
% Reconciled against the runs rather than against memory: the MODFLOW version is
% recorded in the first lines of each scenario's mfsim.lst, and all 46 simulations
% the figures read report 6.7.0. Seven directories under results/gp still record
% 6.5.0 -- two early tracer experiments (_bcfull, _livefull) and the five
% unsuffixed fine-grid runs superseded by the _instbnd sweep -- but none of those
% is read by any figure or statistic in this paper.
The versions used are D-Flow FM kernel 1.2.184, MODFLOW~6 version 6.7.0, and SWMM version 5.2.4 accessed through PySWMM 2.0.1.

\section{Coupling framework}

The coupled simulation is driven by the D-Flow FM model time. D-Flow FM advances in increments of its user time step, $\Delta t_u$, and MODFLOW~6 is advanced at a coupling interval, $\Delta t_c$, that is an integer multiple of $\Delta t_u$. The order of execution and the state passed between the models are shown in Figure~\ref{fig:sequence}.

% jdhughes start here
Two distinct time scales drive the coupling. On every user time step SWMM is advanced first, so that its outfall covers the interval D-Flow FM is about to integrate, and the mean discharge over that step is written into the D-Flow FM source term before D-Flow FM advances. Once per coupling interval, after the last user time step of the interval, the D-Flow FM water levels and depths are mapped to the MODFLOW~6 coastal boundary, MODFLOW~6 is advanced one time step, and the simulated boundary flows are returned to D-Flow FM. Figure~\ref{fig:sequence} is read by where each connector begins. A connector leaves the producing model time step at the point where the exchanged value is formed---the end of that step for a quantity taken instantaneously, the middle for a mean over one step, the whole span for a mean over an interval---and enters the start of the step that consumes it. Its horizontal run is therefore the offset between production and use, measured on the time axis.

The two directions of this exchange carry different lags. The water levels are averaged over the coupling interval, by equation~\ref{eq:tmean} below, and are the boundary condition for the MODFLOW~6 step spanning that same interval, so they introduce no lag. The boundary flows, in contrast, can only be recovered after MODFLOW~6 has finalized that step, so they are held constant over the next coupling interval and lag the hydrodynamics by $\Delta t_c$. That lag is the principal coupling error, and it is what the sweep over coupling intervals reported below is designed to bound. A single MODFLOW~6 step therefore takes its coastal boundary from its own span and its sewer flux from the span before it. The aquifer--sewer seepage is returned on the same schedule as the boundary flows, and so carries the same lag. It is two-way: it is computed from the simulated heads and the water surface in the pipe together, and the same flux is applied to both models with opposite sign, as an inflow to the SWMM junction and an equal withdrawal from the MODFLOW~6 cell, so that mass leaving one model enters the other exactly. Because each exchange uses the state at the end of the preceding operation, the coupling is fully explicit and no iteration is performed within a coupling interval.

\begin{figure}[htbp]
\centering
\widefig{figures/coupling_sequence.pdf}
\caption{Model time steps, order of execution, and state exchanged among the three simulators over two coupling intervals. Bar width is the model time step each simulator integrates: MODFLOW~6 advances the whole coupling interval, $\Delta t_c$, in one step, while SWMM and D-Flow FM tile that span with one bar per user time step, $\Delta t_u$. Lanes are ordered by execution rather than by physical position. Each connector leaves the producing step at the point where the exchanged value is formed and enters the start of the step that consumes it. Repeated user time steps are drawn faded, connector included, and only the leading one is labeled. The exchange break carries no model time, and the faded stub at the right is the interval that follows.}
\label{fig:sequence}
\end{figure}

The explicit scheme is imposed by what the interfaces provide, not chosen for convenience. An implicit coupling requires every simulator to re-evaluate a time step against revised exchange terms, and only one of the three permits it: the MODFLOW application programming interface gives access to the nonlinear solution within a time step, so the groundwater solution can be iterated \citep{hughes2022modflowapi}. The hydrodynamic and sewer interfaces do not. Both provide calls that bracket a time step and advance it, with no facility to re-solve the same step once an exchanged quantity has changed, so a coupled iteration has nothing to iterate on two of its three legs.

The disparity in time step compounds the difficulty. The hydrodynamic model is advanced one user time step of seconds to minutes at a time, the groundwater model on the coupling interval, and the sewer model routes internally on a step of its own, so the models step at rates differing by up to two orders of magnitude. Even were re-evaluation available throughout, defining a converged coupled state across those scales, and deciding which model repeats how many of its own steps per outer iteration, is not a matter of simply looping the exchange. Quantifying how the explicit solution depends on the coupling interval, which the demonstration below does, is the useful thing that can be established without it.

The coupling interval is the principal numerical control on the approach. The coastal boundary of the groundwater model is forced by a semidiurnal tide with a period of approximately 12.42 hours (\si{\hour}), so a coupling interval that is a substantial fraction of that period samples the tidal signal sparsely and aliases it. Coupling intervals of 15 minutes, 30 minutes, and 1, 2, 4, 8, and 24 hours were simulated on each of three D-Flow FM grids to bound this effect.

\section{Coastal exchange}

Water levels simulated by D-Flow FM are applied to the MODFLOW~6 coastal boundary, which is represented by a General-Head Boundary (GHB) package covering the inner bays and a Constant-Head (CHD) package covering the open coast and the lateral perimeter. The D-Flow FM grid and the MODFLOW~6 grid are not coincident, so the exchange is mediated by area-weighted mapping matrices computed once, before simulation, from the intersection of the two grids.

D-Flow FM resolves the tide on a time step orders of magnitude shorter than the coupling interval, so a single boundary value has to summarize everything that happened over the interval. MODFLOW~6 advances that interval in one backward-in-time step, and the exchange term of such a step is applied across its whole length: the volume exchanged is the product of the interval, the conductance, and the head difference. The boundary value must therefore represent the interval rather than any instant within it, and the representation used here is a time average weighted by the wetted state,

\begin{equation}
  \bar{\delta}_{j} = \frac{1}{\Delta t_{c}}
      \int_{t_{k}}^{t_{k+1}} \delta_{j}(t) \, \mathrm{d}t ,
  \qquad
  \bar{s}_{j} = \frac{\int_{t_{k}}^{t_{k+1}} \delta_{j}(t) \, s_{j}(t) \,
      \mathrm{d}t}{\int_{t_{k}}^{t_{k+1}} \delta_{j}(t) \, \mathrm{d}t} ,
  \label{eq:tmean}
\end{equation}

\noindent where $\bar{\delta}_{j}$ is the fraction of the coupling interval over which D-Flow FM cell $j$ is wet (dimensionless), $\delta_{j}(t)$ is unity where that cell is wet and zero where it is dry (dimensionless), $s_{j}(t)$ is the water level it simulates (L), $\bar{s}_{j}$ is the mean of that water level over the wetted part of the interval (L), and $\Delta t_{c}$ is the coupling interval (T). The integrals are evaluated by the trapezoidal rule on the D-Flow FM user time step. Weighting the level by the wetted state is necessary: an unweighted mean would credit an intertidal cell with a water level over the part of the interval it stood dry.

The boundary head assigned to GHB cell $i$ is

\begin{equation}
  h_{i}^{\mathrm{GHB}} = \gamma \sum_{j} w_{ij}^{\mathrm{GHB}} \, \bar{s}_{j} ,
  \label{eq:ghbhead}
\end{equation}

\noindent where $h_{i}^{\mathrm{GHB}}$ is the boundary head of cell $i$ (L), $w_{ij}^{\mathrm{GHB}}$ is the area-weighted contribution of D-Flow FM cell $j$ to MODFLOW~6 cell $i$ (dimensionless), and $\gamma$ is the unit conversion from meters to feet (dimensionless). Constant-head values are computed identically from equation~\ref{eq:ghbhead} using the corresponding CHD mapping.

Intertidal cells alternately wet and dry over a tidal cycle. A D-Flow FM cell with zero water depth contributes no head and no conductance, so the GHB conductance is reduced in proportion to the wetted fraction of the contributing cells,

\begin{equation}
  C_{i}^{\mathrm{GHB}} = C_{i}^{0} \sum_{j} w_{ij}^{\mathrm{GHB}} \,
      \bar{\delta}_{j} ,
  \label{eq:ghbcond}
\end{equation}

\noindent where $C_{i}^{\mathrm{GHB}}$ is the conductance applied to cell $i$ (L$^{2}$T$^{-1}$), and $C_{i}^{0}$ is the conductance of cell $i$ in the base model (L$^{2}$T$^{-1}$). Equations~\ref{eq:tmean} to \ref{eq:ghbcond} together reproduce the interval-mean exchange exactly rather than approximately. Because the instantaneous conductance is the base conductance multiplied by a quantity that is either zero or one, and because MODFLOW~6 carries a single head over the step,

\begin{equation}
  \frac{1}{\Delta t_{c}} \int_{t_{k}}^{t_{k+1}} C^{0} \delta_{j}(t)
      \left( s_{j}(t) - h \right) \mathrm{d}t
    = C^{0} \, \bar{\delta}_{j} \left( \bar{s}_{j} - h \right) ,
  \label{eq:exact}
\end{equation}

\noindent where $h$ is the groundwater head at the end of the step (L). In equation~\ref{eq:exact} the wetted fraction is thus the conductance multiplier and the wet-weighted mean is the boundary head, and no other pair of values reproduces the interval-mean flux. The base-model conductance, rather than the conductance from the preceding coupling interval, is used on the right-hand side of equation~\ref{eq:ghbcond}; applying the multiplier to the previous value would drive the conductance monotonically toward zero over the simulation.

Boundary flows simulated by MODFLOW~6 are returned to D-Flow FM as an external volume source. The source applied to D-Flow FM cell $j$ is

\begin{equation}
  Q_{j}^{\mathrm{ext}} = -\beta \, \delta_{j}
    \left( \sum_{i} a_{ji}^{\mathrm{GHB}} Q_{i}^{\mathrm{GHB}}
         + \sum_{i} a_{ji}^{\mathrm{CHD}} Q_{i}^{\mathrm{CHD}} \right) ,
  \label{eq:qext}
\end{equation}

\noindent where $Q_{j}^{\mathrm{ext}}$ is the volumetric source applied to D-Flow FM cell $j$ (L$^{3}$T$^{-1}$), $a_{ji}$ is the area-weighted contribution of MODFLOW~6 boundary $i$ to D-Flow FM cell $j$ (dimensionless), $Q_{i}$ is the boundary flow simulated by MODFLOW~6 (L$^{3}$T$^{-1}$), and $\beta$ is the unit conversion from cubic feet per day to cubic meters per second (dimensionless). MODFLOW~6 reports boundary flow as positive into the groundwater model, so the leading negative sign renders $Q_{j}^{\mathrm{ext}}$ positive into the surface water. Dry cells receive no source, because $\delta_{j}$ in equation~\ref{eq:qext} is zero where the cell is dry.

\subsection{Representing the tide over the coupling interval}
\label{sec:reduction}

Averaging is one of two representations in common use. The other is to sample the D-Flow FM state at the end of the coupling interval and apply it unchanged, which is the simpler implementation and is consistent in the sense that it converges as the interval is shortened. The two are not better and worse versions of one approximation. They fail by different mechanisms, and which mechanism is active is set by the tide rather than by the model. Sampling preserves the amplitude of the boundary signal but aliases it once the coupling interval exceeds the Nyquist limit of the dominant tidal constituent. Averaging cannot alias, but damps the amplitude by an amount that grows with the interval.

Figure~\ref{fig:reduction} shows what sampling and averaging do to cells of differing wetness. The water level used there is the sum of the three semidiurnal constituents fitted by least squares to the National Oceanic and Atmospheric Administration record at Kings Point, New York (station 8516945, referenced to the North American Vertical Datum of 1988), in mid-Long Island Sound: 1.110 meters (\si{\meter}) for $M_2$, \SI{0.225}{\meter} for $S_2$, and \SI{0.287}{\meter} for $N_2$. The spring--neap and perigean modulation is therefore that of the measured record rather than that of a single sinusoid, which matters because the behavior described below arises from how that record and the coupling interval interact. For a cell that is never dry, the conductance is fixed and only the head is at issue, and the trade between aliasing and damping is the whole of it. For an intertidal cell the conductance carries the more consequential difference. Sampled instantaneously, such a cell is reported as either fully connected or entirely disconnected depending on where the sampling instant happens to fall in the tide, so the boundary switches on and off at the difference between the coupling frequency and the tidal frequency while nothing physical about the cell has changed. The wetted fraction of equation~\ref{eq:tmean} instead reports the partial connection the interval actually contained. What either choice costs in a coupled simulation is reported in Section~\ref{sec:results}.

\begin{figure}[htbp]
\centering
\widefig{figures/boundary_reduction.pdf}
\caption{Representing a tidal boundary by one value per coupling interval, for cells wet 100, 75, 50, and 25 percent of the time, at coupling intervals of 2, 8, and 24 hours. Each pair of panels shows the water level and, beneath it, the multiplier applied to the base conductance. Bed elevations are placed at quantiles of the water level to give the stated wetted fractions. Crosses mark coupling intervals for which instantaneous sampling reports the cell dry, so that no boundary is applied at all. The 2-hour interval is below the $M_2$ Nyquist limit of \SI{6.21}{\hour} and the other two are above it.}
\label{fig:reduction}
\end{figure}

\section{Sewer exchange}

Exchange between the aquifer and the sewer is computed at a set of junctions that fall within the active area of the groundwater model. Each junction is associated with the MODFLOW~6 cell containing it, and a well is added to the groundwater model at that cell.

Infiltration and exfiltration are driven by the difference between the groundwater head and the water surface inside the pipe, and are limited by a conductance representing the leakage area of the pipe wall served by the junction. A per-junction conductance is computed from the sewer geometry as

\begin{equation}
  C_{j} = \lambda \, \pi D_{j} L_{j} ,
  \label{eq:pipecond}
\end{equation}

\noindent where $C_{j}$ is the leakage conductance of junction $j$ (L$^{2}$T$^{-1}$), $\lambda$ is a bulk pipe-wall leakance (T$^{-1}$), $D_{j}$ is the conduit diameter (L), and $L_{j}$ is the length of pipe served by junction $j$ (L), taken as half the length of every conduit that the junction touches. The served length varies by more than a factor of twenty across the network, from approximately 8.5 to \SI{206.5}{\meter}, so a conductance assigned uniformly per connection cannot represent the network geometry.

The driving head depends on whether the water table reaches the pipe. Where the groundwater head stands above the pipe invert the surrounding material is saturated and the gradient is measured to the water surface within the pipe; where the head falls below the invert the pipe free-drains through unsaturated material and the gradient is set by the depth of water in the pipe alone. The driving head is therefore

\begin{equation}
  \Delta h_{j} =
  \begin{cases}
    h_{j} - \left( z_{j} + d_{j} \right), & h_{j} > z_{j} \\[4pt]
    -d_{j},                                & h_{j} \le z_{j}
  \end{cases}
  \label{eq:pot}
\end{equation}

\noindent where $\Delta h_{j}$ is the driving head at junction $j$ (L), $h_{j}$ is the simulated groundwater head in the cell containing junction $j$ (L), $z_{j}$ is the pipe invert elevation (L), and $d_{j}$ is the depth of water in the pipe above the invert (L). Positive $\Delta h_{j}$ produces infiltration into the sewer and negative $\Delta h_{j}$ produces exfiltration. Treating the disconnected case in equation~\ref{eq:pot} as $-d_{j}$ rather than $h_{j} - z_{j}$ is the same treatment given a disconnected stream, and it is necessary: the latter form produces exfiltration from an empty pipe that grows without bound as the water table falls.

The exchange at junction $j$ is

\begin{equation}
  Q_{j} = \eta_{j} C_{j} \Delta h_{j} ,
  \label{eq:swmmq}
\end{equation}

\noindent where $Q_{j}$ is the exchange (L$^{3}$T$^{-1}$), and $\eta_{j}$ is unity for infiltration and an exfiltration ratio less than unity for exfiltration (dimensionless). Exfiltration into saturated material is more difficult than infiltration from it, and the asymmetry is retained. The exchange is applied to SWMM as a generated inflow at the junction and to MODFLOW~6 as an equal and opposite well flux, so that mass is conserved between the two models. Cells that are dry or inactive return a head of $-10^{30}$ in MODFLOW~6 and are assigned zero exchange rather than being treated as real heads.

\section{Sewer discharge tracer transport}

Flow reaching the outfall is introduced to D-Flow FM as a point source carrying a conservative tracer of sewer discharge, which is then advected and dispersed by the hydrodynamic model. The source discharge is the total inflow simulated by SWMM at the outfall, and the tracer concentration assigned to that discharge is constant, so the tracer mass flux follows the simulated SWMM outfall discharge directly.

The SWMM outfall discharge is written into the D-Flow FM source term at each user time step. SWMM is advanced one user time step ahead of D-Flow FM so that its outfall covers the interval D-Flow FM is about to integrate, and the discharge applied over that interval is the mean

\begin{equation}
  \bar{Q}_{s} = \frac{V\!\left(t + \Delta t_{u}\right) - V\!\left(t\right)}
                     {\Delta t_{u}} ,
  \label{eq:blockmean}
\end{equation}

\noindent where $\bar{Q}_{s}$ is the mean SWMM outfall discharge over the user time step (L$^{3}$T$^{-1}$), $V$ is the cumulative volume conveyed through the SWMM outfall (L$^{3}$), and $\Delta t_{u}$ is the D-Flow FM user time step (T). Forming the mean from the cumulative volume, as equation~\ref{eq:blockmean} does, conserves mass exactly over the step, whereas sampling the SWMM outfall discharge instantaneously once per coupling interval and holding it would introduce the same aliasing that the coupling interval itself is being tested for. Advancing SWMM at the user time step rather than at the coupling interval does not alter the sewer solution, because SWMM routes internally at its own routing step in either case and the exchange computed from equation~\ref{eq:swmmq} is still refreshed only once per coupling interval; it raises the resolution at which the outfall is observed.

\section{Implementation}
\label{sec:implementation}

% Trimmed for Advances in Water Resources, whose scope is the water rather than the
% program, and on the house rule that software is the instrument and not the subject.
% What stays here is what changes the water: a source term written in the wrong place
% delivers no volume, and one setting rather than two compiled simulators is what makes
% the comparison controlled. The call sequence each simulator provides is Appendix A,
% where a reader reproducing the work will look for it and a reader following the
% argument will not trip over it.

The coupled simulation is implemented in Python. The three simulators are loaded as shared libraries into a single process, so that state is passed between them as arrays in memory.

Where the sewer source term is written within the D-Flow FM time step decides whether any of it is delivered. D-Flow FM applies its external forcings at the start of a user time step and then computes the flow solution, so a source term written around the complete time step is overwritten by the forcing update before the flow computation reads it, and none of the outfall discharge reaches the receiving water; written between those two operations, all of it does. This is the transfer drawn as $\bar{Q}_s$ in Figure~\ref{fig:sequence}, and the arrow falls within a user time step rather than between two of them. Neither simulator offers that access through the Basic Model Interface itself, and the calls each provides instead are given in \ref{app:interface}, so a coupling of this kind is portable only to the extent that a simulator supplies some equivalent.

\subsection{Coupling strategy}

Because no simulator is modified, what is exchanged, how it is reduced to a single value, and where within a model time step it is written reside entirely in the driver. The cost of trying an exchange formulation is then the cost of running the models rather than of editing and recompiling a hydrodynamic or a groundwater simulator, and two of the choices reported here were settled on that evidence rather than by argument: the representation of the coastal boundary, and the placement of the sewer source term above, which was found by simulating both placements and comparing the volume delivered.

The representation of the D-Flow FM boundary as one value per coupling interval (Section~\ref{sec:reduction}) is selected without altering any of the three models, so the instantaneous and the time-averaged forms were both simulated, at every coupling interval, from the same executables and the same model input. That is what makes the comparison in Figure~\ref{fig:averaging} a controlled one: a single thing differs between the two curves, and the aliasing threshold it exposes is a property of the representation rather than of the program that produced it. Implementing either formulation inside a simulator would have required a separately compiled simulator for each, and would have left the two unavailable simultaneously.

\section{Demonstration}

The coupling is exercised on the Greenport sewer network described above. The groundwater model is a cutout of the regional Long Island aquifer system model \citep{walter2024longisland}, which was developed with MODFLOW~6 to simulate water levels and the position of the saltwater interface from 1900 to 2019. Three D-Flow FM grids of increasing resolution were used---6,491, 16,666, and 41,091 cells---spanning the same domain, sharing a user time step of 300 seconds (\si{\second}), within which the computational step is held by the Courant condition to at most \SI{30}{\second}, and driven by the same tidal and atmospheric forcing over the same 89-day period. The groundwater and sewer models are identical across all simulations, so differences among grids reflect the resolution of the coastal water level alone, and differences among coupling intervals reflect the coupling alone. No metered infiltration record exists for this sewer, so the bulk pipe-wall leakance of equation~\ref{eq:pipecond} is based on published allowances rather than calibrated to measurement: it is set to 0.0655 per day (\si{\per\day}), the value at which the simulated network infiltrates within the range of the design allowances in \cite{asce2007gravity}, and used in every simulation. Each grid was simulated at seven coupling intervals---15 and 30 minutes and 1, 2, 4, 8, and 24 hours---under both sampling and averaging of the coastal boundary, with two further intervals, 6 and 12 hours, added on the coarse grid, giving 46 simulations in total. Every interval divides a day, so each daily stress period holds a whole number of MODFLOW~6 time steps, one to each coupling interval (ninety-six at 15-minute coupling and one at daily), and the coupling steps coincide with the stress period boundaries. The coupling region and the three meshes are shown in Figure~\ref{fig:region}.

\begin{figure}[htbp]
\centering
\widefig{figures/coupling_region.pdf}
\caption{The coupling region at the three hydrodynamic resolutions: (A) coarse, (B) medium, and (C) fine D-Flow FM meshes. The MODFLOW~6 coastal boundary cells and the Greenport sewer network are drawn in every panel because they are identical in every simulation; the groundwater model occupies the area those boundary cells enclose. The meshes extend well beyond the area shown, spanning Long Island Sound, of which the coupling region is a corner. Base map tiles courtesy of the U.S. Geological Survey.}
\label{fig:region}
\end{figure}

\section{Results}
\label{sec:results}

The consequence for the coupled solution was evaluated by comparing each of those simulations against a 15-minute simulation, which resolves the tide and serves as the reference. Both a sampled and an averaged 15-minute simulation were used, because the averaged one shares its formulation with half of the simulations being evaluated; the two give the same result, so the comparison does not depend on that choice.

\subsection{Sensitivity to the representation of the tide}
\label{sec:representation}

The comparison is summarized in Figure~\ref{fig:averaging}. At daily coupling, where the choice between sampling and averaging has consequence, sampling carries 33 millimeters (\si{\milli\meter}) of head error against \SI{3}{\milli\meter} for averaging on the coarse grid. The medium and fine grids give \SI{39}{\milli\meter} against \SI{3}{\milli\meter} and \SI{31}{\milli\meter} against \SI{3}{\milli\meter}, ratios of 12.0, 11.4, and 10.4, so the result does not depend on the discretization. At intervals of 8 hours and shorter sampling and averaging differ by amounts too small to matter (root-mean-square head error stays below \SI{2.2}{\milli\meter} on all three grids), and there sampling is the marginally more accurate of the two at fifteen of the sixteen combinations of grid and interval, which is the damping penalty that averaging pays; that ordering, however, belongs to this start date rather than to the two representations, for the reason established below. Aquifer--sewer seepage behaves as the head does, favoring averaging at daily coupling by factors of 1.4, 2.9, and 2.7.

The sewer discharge tracer is far less sensitive to the choice, though not equally so on every grid. Its 99.99th percentile stays within 1 percent of the reference at every interval on the coarse grid and within a third of a percent on the medium grid, where sampling and averaging cannot be told apart. On the fine grid they can: the sampled percentile sits 3.3 to 4.1 percent from the reference between 1- and 8-hour coupling, whereas the averaged one grows with the interval, from 3.8 percent at 1 hour to 11.1 percent at 8 hours, and at daily coupling the two stand at 18.4 and 11.7 percent. The tracer therefore separates the two representations only where the hydrodynamic grid is fine enough to resolve the plume. At 1 and 2 hours the two departures differ by a few tenths of a percentage point and which of them is nearer the reference carries no weight; from 4 hours on, the tracer orders the representations as the head does.

What distinguishes the two error mechanisms is that only one of them is monotonic in the coupling interval. A truncation error grows smoothly as the step is lengthened, and the averaged solution does exactly that, from \SI{0.05}{\milli\meter} at 30 minutes to \SI{2.7}{\milli\meter} at daily coupling. The sampled solution does not: its head error is \SI{1.3}{\milli\meter} at 8 hours, \SI{36.9}{\milli\meter} at 12 hours, and \SI{32.9}{\milli\meter} at 24 hours, so halving the coupling interval from one day to twelve hours makes the solution worse rather than better. Sampling every 12 and every 24 hours aliases $M_2$ to the same apparent period of 14.8 days, and the two carry comparable error; sampling every 8 hours aliases it to an apparent period of 22.5 hours, which the aquifer damps and which costs almost nothing even though 8 hours is already above the Nyquist limit. It is the period onto which the constituent is aliased, rather than the length of the interval or the crossing of the limit in itself, that sets the error.

The two mechanisms differ in whether they depend on the instant the simulation begins. Averaging flattens the tide within an interval, and how much of it survives is set by the length of that interval against the period of the constituent and by nothing else: 44 percent of the $M_2$ amplitude remains at 8 hours. Sampling takes a water level at an instant and so keeps the whole of the amplitude at every interval. Applied to the fitted water level of Figure~\ref{fig:reduction}, averaging reproduces the mean exactly at every interval and every start while retaining 43 percent of the range at 8 hours, and sampling retains the full range while displacing the mean by up to \SI{6}{\milli\meter} as the start is moved through a day. The aquifer turns the flattened tide into a lowered water table, because a tide holds the mean water table above mean sea level where the flow is unconfined \citep{philip1973periodic, nielsen1990tidal} and a smaller tide holds it lower. That lowering is most of what the error at short intervals is made of: at 8-hour coupling on the medium grid the averaged solution carries \SI{1.5}{\milli\meter} of bias within \SI{2.2}{\milli\meter} of error, against \SI{0.2}{\milli\meter} within \SI{1.0}{\milli\meter} for the sampled one. Averaging carries no start dependence at all: its mean is exact by construction and its damping is fixed by the length of the interval against the period of the constituent. Sampling's displacement does move with the start, although below the Nyquist limit it stays within \SI{3.4}{\milli\meter} of the true mean over a full day of start times on the fitted record, and the aquifer attenuates it: the same 8-hour comparison sets several millimeters at the boundary against \SI{0.2}{\milli\meter} of head bias.

\begin{figure}[htbp]
\centering
\widefig{figures/boundary_averaging.pdf}
\caption{Error in the coupled solution relative to a 15-minute reference simulation, for sampling and averaging of the D-Flow FM boundary, on the coarse grid. (A) simulated groundwater head, (B) aquifer--sewer seepage, (C) sewer discharge tracer, and (D) the 99.99th percentile of the sewer tracer, as a percent departure from the reference shown dashed. The first five days are excluded as model spin-up. The vertical line is the Nyquist limit for the $M_2$ constituent, \SI{6.21}{\hour}. The sampled head error in (A) is not monotonic in the coupling interval: it is larger at 12 hours than at 24, because both intervals alias $M_2$ to the same apparent period of 14.8 days, whereas the 8-hour interval aliases it to an apparent period of 22.5 hours that the aquifer damps. Panel (D) is an upper-tail percentile rather than the maximum of the field, which falls in a single cell and is not a domain-wide measure.}
\label{fig:averaging}
\end{figure}

\subsection{Dependence on the simulation start instant}

That asymmetry constrains the two errors but does not by itself settle which of them is the smaller, and the coupled models were run to find out. Three further start instants were simulated on the coarse grid, at 15:25, 19:20, and 23:20 on the first day, taken from a sweep of the simulated boundary water level as the $M_2$ phases that displace the sampled mean least and most at 6-hour coupling together with one between them. Each was run at 4- and 6-hour coupling, both below the Nyquist limit, under both representations, and each was scored against a 15-minute sampled simulation begun at the same instant, since a simulation can only be compared against a reference that shares its start: fifteen coupled simulations in all. The stop time was held while the start moved later, so the shifted simulations span 88.53 to 88.86 days against the 89.00 of the production start, and errors are compared within a start and never across starts.

The ordering does not survive that test. Root-mean-square head error stays below \SI{1.8}{\milli\meter} at every start and interval, so the two continue to differ by amounts too small to matter, but sampling is the more accurate at only three of the eight combinations of start and interval and averaging at the other five. The margins run from \SI{0.005}{\milli\meter} to \SI{0.737}{\milli\meter}, and the spread across starts, \SI{0.73}{\milli\meter}, is wider than every one of them but the largest. Neither does the ordering follow the displacement of the boundary mean that is supposed to produce it: the start whose sampled boundary mean is displaced most at 6-hour coupling is the one where sampling is furthest ahead in head, and the start displaced least is where it is furthest behind.

What the coupled simulations do confirm is the asymmetry of the mechanisms themselves. The averaged head bias is set by the coupling interval and hardly at all by the start, lying between \SI{-0.45}{\milli\meter} and \SI{-0.48}{\milli\meter} at every 4-hour start and between \SI{-0.70}{\milli\meter} and \SI{-0.84}{\milli\meter} at every 6-hour start, whereas the sampled bias ranges from \SI{1.15}{\milli\meter} to \SI{-1.39}{\milli\meter} and changes sign with the start. Less of the boundary displacement is attenuated than the 8-hour comparison implies: moving the start moves the sampled mean of the simulated boundary water level by \SI{14.6}{\milli\meter} at 6-hour coupling and the sampled head bias by \SI{2.5}{\milli\meter}, a sixth of it rather than a twentieth. Aquifer--sewer seepage is the one comparison that does not turn over: sampling is the more accurate at all eight combinations, by factors of 2.6 to 4.2, the larger at 4-hour coupling and the smaller at 6-hour. Below the Nyquist limit, then, sampling and averaging differ in head by amounts too small to matter, and which of them is ahead depends on the instant the simulation begins.

Where sampling aliases, the start time does matter. The error there depends on which tidal phases the sampling instants occupy, and on the same fitted record the mean of the sampled water level moves by as much as \SI{250}{\milli\meter} at 12- and 24-hour intervals as the start is shifted through a day, against \SI{6}{\milli\meter} at 8 hours; roughly one start in twenty places it within \SI{20}{\milli\meter} of the true mean. The size of the daily-coupling penalty reported above is therefore particular to this start date. What does not depend on the start is that averaging removes the possibility of such a penalty, which is a further reason to prefer it.

Errors of a few tens of millimeters in head are consequential here because they are displacements of the simulated mean rather than scatter about it, and because the questions put to a coastal water table are threshold crossings. Whether the table stands above the pipe invert decides whether a reach infiltrates at all, and whether it stands above a basement floor or a road subgrade decides whether groundwater reaches the surface; groundwater flooding of that kind is assessed by how often and by how far such a level is exceeded \citep{liu2021predicting, su2022slr}. A displacement held for the length of a simulation moves every crossing in the same direction, which is what separates it from the run-to-run variability a modeler is accustomed to discounting.

\subsection{Volume and rate of the coastal exchange}

The comparison to this point is of the aquifer's response to the coastal boundary. The choice between sampling and averaging also sets how much water that boundary delivers, which the response does not report and which is the quantity a water budget is written from. Figure~\ref{fig:volume} shows the cumulative coastal exchange over the 89 days as a percent departure from the 15-minute simulation on the same grid. Under averaging, the volume exchanged is a property of the tide and the aquifer rather than of the coupling: it stays within 2.0 percent of the reference at every interval on all three grids and departs by no more than 1.3 percent at daily coupling. Under sampling it is not. Daily coupling delivers 26.7, 28.5, and 18.3 percent more water to the aquifer on the coarse, medium, and fine grids, and 12-hour coupling on the coarse grid delivers 36.6 percent more, again exceeding the daily interval and for the reason the head error does.

The excess is the intertidal behavior of Figure~\ref{fig:reduction} expressed as a volume. A cell reported wet at the sampling instant is connected at a stage near high water for the whole of the following interval, and one reported dry is disconnected for the whole of it, so at intervals long enough to sample an arbitrary tidal phase the boundary alternately over- and under-delivers, and because the phase sampled drifts slowly rather than randomly the two do not cancel. The wetted fraction and the wet-weighted stage of equation~\ref{eq:tmean} instead deliver the flow the interval contained. The distinction matters more for a water budget than the head errors above do: a coastal exchange inflated by a quarter is a statement about how much water entered the aquifer, not about how accurately its head was resolved. Nor would it be found by shortening the interval and watching the answer settle, because the sampled departure is 1.6 percent or less at every interval at or below the Nyquist limit. Because a coupling interval is normally chosen for computational cost rather than for its relation to the tide, because the penalty cannot be relied upon to fall as the interval is shortened, and because it is carried in the volume exchanged and not only in the accuracy of the response, averaging is preferred here. Every simulation was run under both, and the results that follow are drawn under averaging wherever they do not compare the two.

\begin{figure}[htbp]
\centering
\onecolfig{figures/coastal_reduction.pdf}
\caption{Cumulative coastal exchange over the 89-day simulation, as a percent departure from the 15-minute simulation on the same grid, for sampling and averaging of the D-Flow FM boundary on each of the three hydrodynamic grids. Positive is a larger volume carried into the aquifer. The vertical line is the Nyquist limit for the $M_2$ constituent, \SI{6.21}{\hour}. The 6- and 12-hour intervals were simulated on the coarse grid alone. Averaging holds the exchanged volume within 2.0 percent of the reference at every interval and on every grid, and sampling does not.}
\label{fig:volume}
\end{figure}

The two coupling exchanges respond to the interval differently, and Figure~\ref{fig:coastal} shows why the aquifer head is where the requirement is felt. Over the medium grid the cumulative coastal exchange is almost independent of the interval, the seven simulations spanning 2.0 percent of a total of \SI{281}{\milli\meter} over the 89 days. The rate is not: its root-mean-square falls from 13.9 to 5.7 millimeters per day (\si{\milli\meter\per\day}) between 15-minute and daily coupling, so a long interval retains the volume exchanged and loses most of the tidal signal in it. The medium grid is shown because its interval sensitivity is the largest of the three. That the volume survives the interval is a property of averaging rather than of the coupling: sampled, the same seven simulations span 28.5 percent instead of 2.0 (Figure~\ref{fig:volume}).

Water discharging through the coastal boundary carries dissolved constituents, and what becomes of them is decided in the shallow sediments it passes through, where fresh groundwater mixes with recirculating seawater and the reactions that consume nitrate take place \citep{moore1999subterranean, moore2010effect}. Submarine groundwater discharge is an established nitrogen source to coastal waters, Long Island Sound among them \citep{slomp2004nutrient, tamborski2017submarine}, and the removal along that path is limited by the time water spends in the reactive zone \citep{kroeger2008nitrogen, santos2021submarine}. Tidal exchange is what sets that time: the oscillation drives seawater into the sediment and back out again, and the recirculated volume is commonly a large part of the total discharge \citep{robinson2007effect, santos2012driving, huettel2014benthic}. A simulation that reproduces the volume exchanged over a season while halving the amplitude of the exchange rate, which is what daily coupling does here, is accordingly adequate for a water budget and misleading for a nitrogen one. The coupling interval decides not only how accurately the coupled solution is obtained, but whether the exchange it produces can be asked a reactive question at all.

\begin{figure}[htbp]
\centering
\widefig{figures/coastal_exchange.pdf}
\caption{Coastal exchange against coupling interval on the medium grid, under averaging of the coastal boundary. (A) cumulative exchange over the whole simulation. (B) exchange rate over a ten-day window mid-simulation. Positive is flow into the aquifer. The curves in (A) lie within 2.0 percent of one another while the amplitude in (B) more than halves, so the interval governs the dynamics of the coastal exchange and not the volume of it.}
\label{fig:coastal}
\end{figure}

\subsection{Sewer exchange and hydrodynamic resolution}

The requirement the coupling places on the coupling interval does not depend on the discretization, but the sewer exchange itself does. Figure~\ref{fig:sewer} compares the aquifer--sewer exchange simulated by the three grids at a 15-minute interval, where the coupling contributes least and what remains is hydrodynamic resolution. Net infiltration over the simulation is 16.7, 7.4, and 8.1 liters per day per millimeter of pipe diameter per kilometer of sewer (\si{\liter\per\day\per\milli\meter\per\kilo\meter}) on the coarse, medium, and fine grids. The coarse grid returns more than twice the medium one, and the sequence is not monotonic: refining from medium to fine raises the rate by 10 percent rather than continuing to lower it. A newly constructed gravity sewer is required to leak no more than \SI{9}{\liter\per\day\per\milli\meter\per\kilo\meter} under a hydrostatic test \citep[100 gallons per day per inch of diameter per mile, Paragraph 33.94 of][]{glumrb2014wastewater}, and the infiltration allowances individual agencies adopt for design span roughly an order of magnitude, from 2.3 to \SI{28}{\liter\per\day\per\milli\meter\per\kilo\meter} \citep[Table 3-8 of][]{asce2007gravity}. The coarse grid therefore returns nearly twice the leakage permitted in new construction while the medium and fine grids return somewhat less than it, so the choice of hydrodynamic grid alone decides which side of that limit the simulated network falls on.

The mechanism is in Figure~\ref{fig:sewer}B. What differs between the grids is the fraction of junctions standing below the water table, and so exchanging at all, which averages 23.6, 18.0, and 18.7 percent over the same period and follows the same non-monotonic ordering. Coastal water levels set the head along the seaward boundary, the head sets which junctions are connected, and the connected fraction sets the exchange. The sensitivity is amplified along that path rather than merely transmitted: a 24 percent reduction in connected junctions between the coarse and medium grids produces a 56 percent reduction in net infiltration.

Two consequences follow. The sewer exchange is not converged in hydrodynamic resolution over the range simulated here, so an infiltration rate obtained from a coupled model of this kind should be reported together with the grid that produced it. And the bulk pipe-wall leakance $\lambda$ of equation~\ref{eq:pipecond} is not transferable between grids: it is held at \SI{0.0655}{\per\day} in all three simulations so that the differences are resolution alone, with the consequence that a leakance calibrated to reproduce an observed rate on one grid would be wrong by approximately a factor of two on another.

Whether 244 junctions are more than the exchange requires was tested by halving them. A random subsample of 122 was simulated on the medium grid at a 2-hour interval, with the bulk leakance raised from 0.0655 to \SI{0.1292}{\per\day} so that the total leakage conductance, and with it the pipe area represented, is unchanged. Net infiltration rises by 4 percent, from 7.37 to \SI{7.67}{\liter\per\day\per\milli\meter\per\kilo\meter}, and the connected fraction from 18.0 to 19.6 percent. Simulated head is effectively unaffected: the root-mean-square difference between the two solutions is \SI{0.014}{\milli\meter} and the largest difference anywhere in the domain is \SI{0.036}{\milli\meter}, three orders of magnitude below the error a daily coupling interval introduces on the same grid.

The count is not wholly immaterial. The two exchange series correlate at 0.995 but differ by 12 percent of the root-mean-square of the 244-junction series, so halving the junctions moves where and when the exchange occurs more than it moves how much of it there is. The subsample is a single draw taken with a fixed seed, and a comparatively lumped one: its junctions fall in 60 distinct groundwater cells against a mean of 63.6 over 200 draws, so 4 percent is nearer an upper bound on the effect than a typical value.

What the result does not license is reducing the junction count without compensating for it. Conductance is assigned per junction from the pipe area that junction serves, so a subsample carries only the area it happens to include; at the unchanged leakance these 122 junctions represent slightly over half the conductance of the full set. What the exchange follows is the total leakage area represented, not the number of points that area is distributed over.

Groundwater entering a sewer occupies conveyance the network would otherwise have for storm flow, and where the water table is tidally forced the two arrive together; the same infiltration is what sewer rehabilitation is undertaken to reduce, and what decides whether a low-lying system floods from below rather than from the surface \citep{liu2018gwsewer, su2019infiltration, liu2021predicting, su2022slr}. A factor of two in simulated infiltration is a factor of two in the quantity such an assessment turns on, and here it is set by a discretization chosen for the coast rather than by anything about the sewer.

\begin{figure}[htbp]
\centering
\widefig{figures/sewer_exchange.pdf}
\caption{Aquifer--sewer exchange simulated by the three D-Flow FM grids at a 15-minute coupling interval under the averaged coastal boundary. (A) cumulative net exchange into the sewer, annotated with the infiltration rate each grid settles at, in liters per day per millimeter of pipe diameter per kilometer of sewer; positive is infiltration. (B) percentage of the 244 coupled junctions standing below the water table and therefore exchanging, plotted at every coupling step. The first five days, shaded, are excluded from the reported rates as model spin-up. The MODFLOW~6 and SWMM models are identical in the three simulations and the bulk pipe-wall leakance is the same, so the differences are those of hydrodynamic resolution alone.}
\label{fig:sewer}
\end{figure}

\subsection{Cost of the coupling interval}

What the coupling interval costs is shown in Figure~\ref{fig:cost}. The marginal cost of one coupling step is \SI{0.25}{\second} and does not depend on the grid, which is what should be expected: the MODFLOW~6 and SWMM models are identical across the three D-Flow FM discretizations, so a shorter interval buys additional solves of the same two models. The hydrodynamic cost it is set against is not constant. D-Flow FM alone accounts for approximately 36, 94, and 239 minutes of the 89-day simulation on the coarse, medium, and fine grids, which is very nearly linear in cell count at 5.5 to $5.9\times10^{-3}$ minutes per cell. The consequence is that the penalty for frequent coupling falls as the hydrodynamic model is refined: increasing the number of exchanges ninety-six-fold, from daily to 15-minute coupling, costs of order 100 percent of the run on the coarse grid but between 13 and 39 percent on the medium and fine grids.

The error incurred by a long coupling interval is not monotonic in the interval, so it cannot be bounded by shortening the interval a little; the cost of a short one is modest, and becomes more modest as the hydrodynamic grid approaches the resolution at which such models are actually run. There is correspondingly little to recommend a long coupling interval on a production-scale grid, and the interval is better chosen from the tidal constituents than from the run-time budget.

Run times were measured on a single workstation with the simulations executed sequentially. Run-to-run variation is a few percent, which is the reason several medium-grid points in Figure~\ref{fig:cost}A fall marginally below their own daily-coupling baseline, and the reason the spread between the two representations at a given grid should be read as scatter rather than as a property of either.

\begin{figure}[htbp]
\centering
\widefig{figures/coupling_cost.pdf}
\caption{Cost of the coupling interval, for both representations of the D-Flow FM boundary on the three grids. (A) run time as a multiple of that series' own daily-coupling simulation. (B) additional run time against additional coupling steps, each series referred to its own daily-coupling simulation rather than to a fitted intercept, so that the collapse shown is a property of the measurements; the line is the pooled marginal cost through the origin. The vertical line in (A) is the Nyquist limit for the $M_2$ constituent, \SI{6.21}{\hour}. Coupling more frequently costs approximately the same number of seconds per exchange on every grid, because the MODFLOW~6 and SWMM models being solved are identical; what differs is the hydrodynamic cost that expense is measured against.}
\label{fig:cost}
\end{figure}

\section{Choosing the coupling for a problem of this kind}
\label{sec:choosing}

Taken together these results prescribe how a coupling of this kind should be set up, without repeating the sweep: (1) identify the dominant tidal constituent and compute, for each candidate interval, the period onto which sampling would alias it, which is the reciprocal of the offset between the constituent frequency and the nearest multiple of the sampling frequency; (2) reject the intervals whose alias period the aquifer cannot damp; (3) average the boundary over the interval rather than sampling it, whatever interval is chosen; and (4) shorten the interval until the solution stops moving.

The second of those, not the Nyquist limit, is the test that matters: at this site 8-hour coupling aliases $M_2$ to \SI{22.5}{\hour}, a period the aquifer damps, whereas 12- and 24-hour coupling alias it to 14.8 days, which it does not. The coastal coupling of \cite{jin2025gwocean}, exchanging every 10 minutes, sits far inside that bound. Coupling error has been analyzed before in integrated surface--subsurface models, where it grows with the rate of change of the exchanged flux and is reduced by adapting the length of the step \citep{fiorentini2015coupling}. A periodic boundary is not covered by that remedy: the error here is set by the period onto which the constituent aliases rather than by any gradient, and it is not monotonic in the interval, so a heuristic that shortens the step when coupling is strong can shorten it from 24 hours to 12 and make the solution worse.

Averaging is preferred at whatever interval is chosen because it cannot alias, because it holds the exchanged volume to within 2 percent at every interval and on every grid simulated here, and because its bias is set by the length of the interval rather than by the instant the simulation begins, which the sampled bias is not (Section~\ref{sec:representation}).

Shortening then buys accuracy at a known price. On these grids the solution stops moving at 2 hours and shorter, where the averaged solution lies within \SI{0.32}{\milli\meter} in head and 0.13 percent in exchanged volume of the 15-minute reference, and each halving of the interval costs one further solve of the groundwater and sewer models per exchange, \SI{0.25}{\second} apiece.

\section{Summary and conclusions}

An approach has been described for coupling a coastal hydrodynamic model, a groundwater flow model, and a sewer network model in memory through their library interfaces, together with the exchange formulation between each pair. None of the three simulators is modified. Each is loaded as a shared library, advanced through the calls its interface provides, and passed state as arrays by a controlling program in which the whole of the coupling resides. The coupling itself is inexpensive, and its cost does not depend on the hydrodynamic discretization: one exchange costs \SI{0.25}{\second} on all three grids simulated, because the groundwater and sewer models being solved are identical in each. What differs is the hydrodynamic cost that expense is set against, which is close to linear in cell count.

The coupling interval introduces an error whose character is set by the tide rather than by the numerics. Because MODFLOW~6 integrates a whole coupling interval in one step, the coastal boundary value must represent the interval and not the instant at its end, and the two representations fail by opposite mechanisms: sampling preserves amplitude but aliases, averaging cannot alias but damps. At intervals of 8 hours and shorter the difference is immaterial, root-mean-square head error remaining below \SI{2.2}{\milli\meter} on every grid. At daily coupling sampling carries 12.0, 11.4, and 10.4 times the error of averaging on the coarse, medium, and fine grids. That the same factor is recovered on three independent discretizations is the strongest evidence available here that the requirement belongs to the coupling rather than to the model.

What separates the two mechanisms is whether the error grows monotonically with the interval. The averaged solution deteriorates smoothly, as a truncation error does; the sampled solution carries more head error at 12 hours than at 24, so shortening the interval can make the solution worse rather than better. Sampling every 12 and every 24 hours aliases the $M_2$ constituent to the same apparent period of 14.8 days and the two carry comparable error, whereas 8-hour sampling aliases it to an apparent period of 22.5 hours that the aquifer damps and that costs almost nothing, although 8 hours already exceeds the Nyquist limit. It is the period onto which the dominant constituent is aliased, rather than the length of the interval or the crossing of the limit in itself, that sets the penalty. A coupling interval is therefore better chosen from the tidal constituents present than from the run-time budget, and a shorter interval cannot be assumed to be a safer one.

Set against that, exchanging often is affordable, and becomes more so as the hydrodynamic model is refined. Increasing the number of exchanges ninety-six-fold costs of order 100 percent of the run on the coarse grid but between 13 and 39 percent on the medium and fine grids (about twice the daily-coupling run, and 1.13 to 1.39 times it), because the fixed hydrodynamic cost grows with resolution while the marginal cost of an exchange does not. At the resolutions such models are ordinarily run at there is accordingly little to recommend a long coupling interval, and none to prefer sampling over averaging.

What the demonstration establishes about the sewer problem is more equivocal. The requirement the coupling places on the interval is grid-independent; the quantity of interest is not. Net infiltration is 16.7, 7.4, and \SI{8.1}{\liter\per\day\per\milli\meter\per\kilo\meter} on the three grids, so the hydrodynamic discretization moves it by more than a factor of two and does so without settling, and the bulk pipe-wall leakance is correspondingly not transferable between grids. Against that, the number of junctions at which the exchange is computed proves nearly immaterial once the total leakage area is held fixed: halving them changes net infiltration by 4 percent and simulated head by \SI{0.014}{\milli\meter}. An infiltration rate obtained from a coupled model of this kind should therefore be reported together with the hydrodynamic grid that produced it, and the discretization of the sewer is the lesser concern.

The coupling depends on one capability the Basic Model Interface does not define. Initializing a simulation, advancing it by a time step, finalizing it, and reading and writing named state are sufficient for the coastal and sewer exchanges, and where a simulator implements the standard the coupling is written directly against it. Returning the sewer outfall discharge to the hydrodynamic model is not covered by those four: the source term must be written after the external forcing update and before the flow computation, and writing it around the combined time-stepping call delivers none of the volume. Both simulators written to within a time step here provide the calls required, by different and non-standard routes: MODFLOW~6 through an interface that extends the standard, and D-Flow FM by providing the phases of its user time step separately. A simulator offering only the standard interface could not be coupled in this way.

Two directions follow. The first is an implicit coupling, in which the exchange is iterated to convergence within a step rather than applied once. That is unavailable today for reasons that lie with the simulators and not with the approach: the hydrodynamic and sewer interfaces advance a time step but cannot re-solve one, so two of the three legs offer nothing to iterate upon, and the three codes step at rates differing by up to two orders of magnitude, which leaves the definition of a converged coupled state unsettled. Establishing how the explicit solution depends on the coupling interval, as is done here, is what can be established without it.

The second is that nothing in the approach is particular to these three codes. The coupling resides in the driver, the simulators are used as distributed, and a different combination requires only that each provide the four capabilities and, for an exchange that must occur within a time step, an equivalent sub-time-step call. Standardizing that call would make this class of coupling portable rather than bespoke, and is the single change to the interfaces that would most extend what can be coupled in memory.

% AdWR asks for no section of this kind, and this one was written to the EMS
% software-availability checklist. What that checklist demanded and nothing else
% wanted -- program language, cost, hardware -- has been cut. What remains is kept
% deliberately: the Option C obligation to cite and link the data deposit is
% answered here in the body rather than only in the submission system, and the USGS
% Fundamental Science Practices requirement that generative-AI use in producing
% CODE be disclosed needs the sentence in the second paragraph, which the AI
% declaration below then matches. Versions are here because a coupled result is not
% reproducible without them. The archive DOI and the data deposit are still to be
% supplied.
\appendix

\section{Interface calls the coupling requires}
\label{app:interface}

The coupling requires each simulator to initialize a simulation from its own input, advance it by a time step, finalize it, and read and write named state. The Basic Model Interface \citep{hutton2020bmi} codifies those four capabilities, and where a simulator implements it the driver calls them by their standard names. MODFLOW~6 and D-Flow FM do; SWMM does not, and is driven through the SWMM toolkit interface, which supplies the same capabilities under different names.

Returning the sewer outfall discharge to the hydrodynamic model requires one capability beyond those four, namely writing model state part way through a time step, and the Basic Model Interface does not define it. The two simulators written to in that way here provide it by different and non-standard routes. MODFLOW~6 extends the Basic Model Interface with functions that give access to the time step and to the solution within it \citep{hughes2022modflowapi}, and the coastal exchange is written through those. D-Flow FM splits its user time step into separate initialization, computation, and finalization calls alongside the combined update, and the tracer source term is written between the first and the second, for the reason given in Section~\ref{sec:implementation}.

Both routes were established against the distributed libraries, requiring no knowledge of the internals beyond the names of the calls and no ability to recompile either simulator. A strategy that needs state a simulator does not expose, or needs it at a moment the simulator does not permit, cannot be written in the driver at all, which is the bound on what this approach can reach.

\section*{Data and code availability}

The three simulators are used as distributed and are not modified by this work. MODFLOW~6 is developed by the U.S. Geological Survey and is available from \url{https://github.com/MODFLOW-ORG/modflow6}; version 6.7.0 was used, with the application programming interface and the \texttt{modflowapi} Python interface. D-Flow FM is developed by Deltares and is available from \url{https://oss.deltares.nl/web/delft3dfm}. SWMM is developed by the U.S. Environmental Protection Agency and is available from \url{https://github.com/USEPA/Stormwater-Management-Model}; version 5.2.4 was used, driven through the \texttt{PySWMM} Python interface, version 2.0.1. All three are loaded as shared libraries by the coupling.

The coupling itself is Python and is the only software new to this work. It comprises the driver that advances the three simulators and moves state between them, and the analysis and figure-generation code from which every figure in this paper is drawn; the latter rebuilds each figure from archived summaries without the simulation output, so the figures can be reproduced from the software release alone. Initial drafts of that code were produced with a generative artificial intelligence tool and were subsequently reviewed and tested by the authors, as described below. % TODO: archive the coupling code and cite the DOI here rather than a bare repository URL.

The model input and the simulation output are published by the U.S. Geological Survey as a data release on ScienceBase, and are cited there rather than reproduced here. % TODO: replace with the ScienceBase DOI and a \citep of the data release once it is published. AdWR is Option C for research data as EMS was, so the deposit must be cited and linked rather than offered on request.

% Sizes, for the data release description: the simulation output is approximately
% 240 GB retained across 46 scenarios, of which the D-Flow FM map files are excluded
% -- they are deleted after each run once the tracer has been verified, and are
% reproducible from the model input. The archived summaries the figures are drawn
% from are five netCDF files totaling under 1 MB and are included in the coupling
% code repository rather than the data release.

\section*{CRediT authorship contribution statement}

% TODO: required by the guide, and only the co-authors can supply it. Roles are
% drawn from the CRediT taxonomy: conceptualization, data curation, formal analysis,
% funding acquisition, investigation, methodology, project administration,
% resources, software, supervision, validation, visualization, writing -- original
% draft, writing -- review and editing.

\section*{Declaration of generative AI and AI-assisted technologies in the manuscript preparation process}

During the preparation of this work the authors used Claude (Anthropic) to produce initial drafts of the manuscript text and of the analysis and figure-generation code in the accompanying software release, and to perform initial analysis of the simulation output. All simulation results were produced by MODFLOW~6, D-Flow FM, and SWMM. After using this tool, the authors reviewed, tested, and edited the content as needed and take full responsibility for the content of the published article.

% Scope is deliberately broader than Elsevier's template, which is written for a tool
% used to improve readability and language. USGS Fundamental Science Practices
% additionally require that the use of generative AI in producing CODE be disclosed,
% and that AI use be described in the information product; the software statement
% above carries the matching sentence. See the FSP FAQ on restrictions on using
% generative AI, and SM 502.10, which reserves authorship to human researchers -- no
% AI tool is or could be listed as an author here.
%
% The sentence naming the three simulators is not filler. It states that no quantity
% reported in this paper originated in a language model: every number comes from the
% physics codes, through analysis code that is archived and re-runnable, and every
% figure rebuilds from committed archives.
%
% NOT disclosed in figure captions, deliberately, and the AdWR guide is explicit
% enough to settle it. Caption-level disclosure is asked of images CREATED with AI
% tools; the same section permits AI tools in data visualizations "when the output
% is directly derived from underlying data using reproducible analytical,
% computational, or statistical methods". All eight figures are that case, as is
% the graphical abstract: scripts in the software release, drawing committed
% archives, rebuilding byte-identically without the simulation output. A caption
% note would imply generated imagery, which would be false.
%
% STILL TO CONFIRM, and not a drafting matter: that the tool used is approved for use
% through the USGS Associate Chief Information Officer and Department of the Interior
% process, which FSP requires before deployment rather than at disclosure.

% Required: "Authors must disclose any funding sources who provided financial
% support for the conduct of the research and/or preparation of the article", with
% the role of the sponsor stated. The guide's format is
%   Funding: This work was supported by the [AGENCY] [grant number xxxx].
% and, if there was none, the sentence to use verbatim is
%   This research did not receive any specific grant from funding agencies in the
%   public, commercial, or not-for-profit sectors.
% Only the authors can say which applies, so the placeholder is left visible in the
% typeset draft rather than hidden in this comment.
\section*{Funding}

This work was conducted by the U.S. Geological Survey as a component of a regional assessment of compound flood hazard from the combined effects of coastal, stormwater, and groundwater emergence flooding \citep{usgs2021compound}. \textbf{[TODO: name the specific program or cooperative agreement supporting that assessment, if any; only the authors can say.]}

% The three closers below are the standard set and the standard order:
% Acknowledgments, then Authors' Note, then Disclaimer. The competing-interest
% sentence lives under Authors' Note and not under Elsevier's own heading, which is
% a deliberate house choice; the declarations tool upload at submission is separate
% and still required either way. The Disclaimer is the bureau's requirement rather
% than the journal's and stays whatever else moves.
%
% KNOWN DEVIATION, deliberate: AdWR asks that acknowledgements sit directly before
% the reference list. Authors' Note and Disclaimer follow them here, because the
% closer order is the house convention and those two are not the journal's to move.
\section*{Acknowledgments}

% Placeholder.

\section*{Authors' Note}

The authors do not have any conflicts of interest to report.

\section*{Disclaimer}

Any use of trade, firm, or product names is for descriptive purposes only and does not imply endorsement by the U.S. Government.

\bibliographystyle{elsarticle-harv}
\bibliography{Hughesetal_AWR_LISSCoupling}

\end{document}
